\clearpage
\section*{Task 2: Autoformer forecasting challenge}
\subsection*{Problem, data, and modeling evidence}
The observed target contains 43,656 regularly spaced samples. The final forecast covers exactly indices 43,657--43,824, a 168-step horizon. The ten aligned external measurements contain six continuous variables (A--F) and four mutually exclusive indicators (G--J). Both historical and future-horizon external values are supplied; future target values remain hidden. Structural checks confirm index continuity, nonmissing nonnegative training targets, hidden test targets, full external coverage, and one-hot indicators. Task 2 errors are in the anonymized target's original units, not Task 1's vibration units; sMAPE is a percentage.

Diagnostics present in the notebook were checked against the supplied data. The target ranges from 0 to 994, with mean 98.17, median 73, standard deviation 90.83, skewness 1.83, and excess kurtosis 5.00. The 99th percentile is 418. The correlation between 168-step rolling mean and standard deviation is 0.787 for raw values and $-0.180$ after $\log(1+y)$, supporting level-dependent variance and a scale-stabilizing transform. This motivates testing a log transform; it does not establish that it will optimize original-scale RMSE.

Global ACF is 0.966 at lag 1, 0.401 at lag 24, 0.158 at lag 48, and 0.028 at lag 168. Local lag-24 ACF over 672-point windows varies from $-0.026$ to 0.637. The periodogram and local ACF were inspected rather than treating one spectral peak as a fixed operating period. The evidence supports strong short-lived dependence and changing local recurrence (Figure~\ref{fig:t2acf}), so copying Task 1's fixed pace routing would lack the same justification.

External associations are modest: contemporaneous correlations are 0.169 for A and $-0.245$ for D; they are diagnostic associations rather than causal claims. The supplied future regime profile differs from history: H occurs in 51.79\% of future steps versus 32.21\% historically, while I falls to 8.93\% from 34.99\%. This supports testing known-future covariates and also cautions against assuming the final horizon has the historical average regime mix.

\subsection*{Chronological validation and preprocessing}
Screening uses seed 2026 and origins 43,320 and 43,488, at most four epochs, patience one, and training-window stride eight. Final evaluation uses seeds 2026 and 3407 and origins 43,152, 43,320, and 43,488. Their nonoverlapping targets are 43,153--43,320, 43,321--43,488, and 43,489--43,656. Each origin permits target history only through that origin. The final stage uses stride four, at most ten epochs, patience two, and 10,663/10,705/10,747 training windows at the three origins. Shuffling eligible training windows does not randomize the temporal holdout.

Each fold fits target location/scale and continuous-covariate scalers solely on pre-origin history. Indicators remain binary. Known-future covariates enter the decoder for the validation horizon; future target values never enter it. Metrics are computed after inversion to original target units and clipping forecasts to nonnegative values. Two seeds and three folds expose initialization and temporal variation at manageable cost; these six evaluations are not six independent datasets.

\subsection*{Autoformer implementation and selection}
The self-contained encoder--decoder follows the decomposition and autocorrelation ideas of Autoformer~\cite{autoformer}, extending the mechanisms explored in Task 1. Moving averages with replicated endpoints split hidden sequences repeatedly after encoder mixing/feed-forward sublayers and decoder self/cross mixing/feed-forward sublayers. Decoder trend contributions are projected and accumulated rather than discarded. Seasonal layer normalization removes the temporal mean.

The mixer takes FFTs of learned queries and keys, forms circular delay correlations, averages scores across heads/features, selects approximately $3\ln L$ delays per example, and softmax-weights shifted values. It therefore aggregates temporal displacements rather than constructing vanilla pairwise dot-product Attention. Encoder inputs contain target history plus projected historical covariates. The decoder starts with 84 historical seasonal/trend positions, zero future seasonal targets, and a history-mean future trend; it also receives the known future covariates. No future target or inferred calendar identity is used.

The question-driven search first tests target representation, then history/decomposition, then capacity; all screening comparisons use the same seed/folds (Table~\ref{tab:t2screen}). Raw global standardization wins mean RMSE 50.49 versus 61.58 for log1p, although log1p has lower sMAPE (55.05 versus 83.49). This demonstrates metric disagreement: stabilizing relative scale need not improve squared error on large target levels. Context 336/kernel 25 wins the structure stage; extending to 672/kernel 49 does not help these folds. Width 96 then improves screening RMSE from 50.49 to 48.36 while increasing parameters from 117,889 to 262,849. These are screening observations, not multiple-seed architecture superiority claims.

\begin{table}[htbp]
\centering\small
\caption{Screening means over one seed and two chronological horizons. Candidate names encode context length, kernel, or model width.}\label{tab:t2screen}
\begin{tabular}{llrrrr}
\toprule
Stage & Candidate & RMSE & MAE & sMAPE (\%) & Parameters \\
\midrule
Target & T\_raw & 50.49 & 36.40 & 83.49 & 117,889 \\
Target & T\_log1p & 61.58 & 36.07 & 55.05 & 117,889 \\
Structure & S\_168\_k25 & 54.37 & 38.35 & 94.66 & 117,889 \\
Structure & S\_336\_k25 & 50.49 & 36.40 & 83.49 & 117,889 \\
Structure & S\_336\_k49 & 52.43 & 36.51 & 96.10 & 117,889 \\
Structure & S\_672\_k49 & 55.98 & 43.72 & 86.09 & 117,889 \\
Capacity & C\_compact32 & 52.76 & 36.93 & 79.64 & 30,273 \\
Capacity & C\_base64 & 50.49 & 36.40 & 83.49 & 117,889 \\
Capacity & C\_medium96 & 48.36 & 33.51 & 84.22 & 262,849 \\
\bottomrule
\end{tabular}
\end{table}

Selection minimizes mean fold RMSE. Candidates within 0.5\% of the best are ordered by parameter count and then mean best epoch before RMSE. The unknown leaderboard penalty coefficients are not estimated, and no leaderboard score is used as a validation signal.

\subsection*{Required external-data ablation and final results}
The final comparison holds the selected architecture and optimization settings fixed, changing only external inputs. Both variants use context 336, label length 84, horizon 168, model width 96, four heads, two encoder layers, one decoder layer, feed-forward width 192, moving average 25, factor 3, and dropout 0.1. AdamW uses learning rate and weight decay $10^{-4}$, batch size 32, standardized raw-target MSE loss, and gradient clipping at 1.0.

\begin{table}[htbp]
\centering\small
\caption{Final ablation: arithmetic means over two seeds $\times$ three folds; RMSE spread is the sample standard deviation across the six run-level RMSEs.}\label{tab:t2ablation}
\begin{tabular}{lrrrrr}
\toprule
External inputs & Mean RMSE $\pm$ SD & MAE & sMAPE (\%) & Parameters & Min/run\\
\midrule
A--J & $58.72\pm16.23$ & 45.25 & 82.61 & 262,849 & 27.65\\
None & $88.41\pm23.55$ & 76.54 & 102.41 & 260,929 & 25.96\\
\bottomrule
\end{tabular}
\end{table}

External inputs reduce mean RMSE by 29.69 (33.58\%), MAE by 31.29 (40.89\%), and sMAPE by 19.81 percentage points. Their extra 1,920 parameters are a 0.74\% increase. Every matched fold/seed RMSE improves (Table~\ref{tab:t2folds}), supporting selection of \texttt{FINAL\_with\_external} under the notebook's rule. This is the measured benefit of all A--J jointly; the experiment does not identify the individual contribution of each feature or separate historical from future-covariate benefit.

\begin{table}[htbp]
\centering\small
\caption{All six final fold/seed pairs. Epoch is the selected checkpoint epoch with external data; the complete per-variant metrics and epoch fields are retained in \texttt{final\_cv\_results.csv}.}\label{tab:t2folds}
\begin{tabular}{rrrrrrr}
\toprule
Origin & Seed & Epoch & \shortstack{RMSE\\with} & \shortstack{MAE\\with} & \shortstack{sMAPE (\%)\\with} & \shortstack{RMSE\\without}\\
\midrule
43152 & 2026 & 4 & 71.11 & 54.76 & 74.26 & 117.06 \\
43152 & 3407 & 4 & 77.09 & 63.29 & 78.22 & 119.52 \\
43320 & 2026 & 2 & 60.27 & 43.98 & 81.40 & 76.25 \\
43320 & 3407 & 1 & 65.73 & 54.44 & 94.92 & 79.93 \\
43488 & 2026 & 2 & 38.48 & 29.47 & 89.28 & 68.64 \\
43488 & 3407 & 5 & 39.62 & 25.54 & 77.56 & 69.04 \\
\bottomrule
\end{tabular}
\end{table}

Temporal difficulty dominates much of the spread. With-external fold-mean RMSEs are 74.10, 63.00, and 39.05 in chronological order; seed-mean RMSEs are 56.62 and 60.81. This favors reporting all folds rather than presenting the easiest final horizon or one seed as overall performance.

Horizon curves pool the six squared errors at each step, then take a square root (Figure~\ref{fig:t2horizon}); the pooled all-step RMSE is 60.56 with external inputs versus 90.98 without, distinct from the arithmetic mean run RMSE in Table~\ref{tab:t2ablation}. With-external block RMSE rises from 38.64 over steps 1--24 to 64.01 over steps 145--168 and reaches 74.97 at steps 121--144. Neither curve increases monotonically, and the external advantage is not uniform at every step.

\subsection*{Full-history training, predictions, and limits}
The final model is freshly initialized with seed 2026 and fitted on all 43,656 observed targets. Its training-window stride is two, with scalers fitted on the full observed history. The six with-external best epochs are 4, 2, 2, 4, 1, and 5; their median selects \textbf{three full-refit epochs}. The checkpoint confirms \textbf{262,849 trainable parameters}; refitting took 34.71 minutes. The saved prediction array and CSV contain exactly 168 finite nonnegative values in chronological order. They are preserved without regeneration; 24 values are zero after clipping. The original forecast is shown in Figure~\ref{fig:t2forecast}. Hidden-target performance is unavailable, and no leaderboard submission was made during preparation.

The epoch ledger distinguishes the final three epochs from preceding training. The original early-stop loops imply 55 screening epochs and 58 final-CV epochs (30 with and 28 without external data), plus one smoke epoch and three refit epochs: \textbf{117 epochs for the complete pipeline}. The prepared metadata uses this conservative whole-pipeline count rather than silently treating three refit epochs as the entire cost. Counts are derived from saved best epochs, patience, and epoch caps; the ledger states that derivation explicitly.

The original notebook does not call \texttt{model.eval()} in validation or final prediction. Consequently, dropout remains active even under \texttt{no\_grad}; reported metrics and forecasts describe that stochastic inference protocol. Also, the same fold targets determine checkpoints and compare architectures, so these are model-selection estimates rather than an independent final test. The three recent folds, changing local recurrence, and shifted future covariate mix limit generalization claims. No expensive rerun or silent replacement of results was performed. The observed external-data advantage supports this selected run, while deterministic evaluation and untouched additional origins would strengthen a future confirmation.

As in Task 1, useful model structure must match the data. Here, local decomposition/delay mixing is paired with legitimately known future measurements, while transform, context, and capacity are chosen by chronological evidence rather than copied from the synthetic benchmark.

\evidencefigure{Question 2 - Leaderboard/figures/global_acf.png}{Task 2: original Kaggle global ACF diagnostic. Short-lag dependence decays substantially before the full 168-step horizon.}{fig:t2acf}
\evidencefigure{Question 2 - Leaderboard/figures/horizon_comparison.pdf}{Task 2: horizon RMSE derived from the saved \texttt{final\_horizon\_rmse.csv}, using a shared vertical scale for both final ablation variants.}{fig:t2horizon}
\evidencefigure{Question 2 - Leaderboard/figures/final_forecast.png}{Task 2: original Kaggle full-history forecast. The continuation has no revealed ground truth; it is not a test-accuracy plot.}{fig:t2forecast}
